% ==============================================================================
% SECTION 3: THEMATIC SYNTHESIS
% ==============================================================================
% Extensión: 2000–3000 palabras
% Estructura: tabla de extracción (longtable), subsecciones por tema,
% ecuaciones display (equation, align, cases), cajas destacadas (infobox, keypoint),
% listings de código (Python, R, bash), tablas comparativas (tabularx)

\section{Síntesis temática}\label{sec:synthesis}

% --- RESUMEN GENERAL ---
Los 15 estudios seleccionados fueron analizados y categorizados en cuatro ejes
temáticos que abordan desde el diseño estructural hasta la validación científica
de los sistemas. La \cref{subsec:architectures} examina las arquitecturas
agénticas; la \cref{subsec:pipeline-automation} analiza la automatización de
pipelines específicos, incluyendo genómica, scRNA-seq y farmacometría; la
\cref{subsec:reproducibility} evalúa las estrategias de reproducibilidad y
benchmarking; y finalmente, la \cref{subsec:governance} discute las limitaciones
técnicas y los marcos de gobernanza.

La \cref{tab:extraction} sintetiza las características operativas, aplicaciones
clave y hallazgos principales de la evidencia analizada, sirviendo como base
para la discusión técnica subsiguiente.

\begin{enumerate}
  \item Arquitecturas agénticas y patrones de orquestación.
  \item Automatización de pipelines bioinformáticos específicos.
  \item Reproducibilidad, benchmarking y evaluación comparativa.
  \item Limitaciones técnicas, supervisión humana y gobernanza.
\end{enumerate}


% --- TABLA DE EXTRACCIÓN (OBLIGATORIA) ---
% Longtable permite que la tabla continúe en varias páginas
\begin{longtable}{p{3.0cm} p{2.8cm} p{2.7cm} p{3.5cm} p{2.6cm}}
\caption{Tabla de extracción de datos de los estudios incluidos.}
\label{tab:extraction}\\
\toprule
\textbf{Autor, año} & \textbf{Enfoque} & \textbf{Aplicación} &
\textbf{Hallazgo clave} & \textbf{Limitación} \\
\midrule\endfirsthead
\multicolumn{5}{l}{\small\itshape Continuación de la página anterior}\\
\toprule
\textbf{Autor, año} & \textbf{Enfoque} & \textbf{Aplicación} &
\textbf{Hallazgo clave} & \textbf{Limitación} \\
\midrule\endhead
\midrule\multicolumn{5}{r}{\small\itshape Continúa en la página siguiente}\\\endfoot
\bottomrule\endlastfoot

\textit{Artificial Intelligence in Transcriptomics}, 2026 &
Modelo IURA &
Transcriptómica &
Transición hacia sistemas autónomos &
Marco emergente \\
\addlinespace
\textit{BioMedAgent}, 2026 &
Framework multi-agente &
Workflows bioinformáticos &
Selección autónoma y encadenamiento de tareas &
Alto coste computacional \\
\addlinespace
\textit{PKGPT}, 2026 &
Agente recursivo basado en LLMs &
Modelado farmacocinético &
Refinamiento iterativo mediante software externo &
Requiere supervisión experta para plausibilidad e interpretabilidad \\
\addlinespace
\textit{ChatSpatial}, 2026 &
Orquestación basada en esquemas &
Transcriptómica espacial &
Reproducibilidad entre workflows R/Python &
Preprint; validación todavía emergente \\
\addlinespace
\textit{CellVoyager}, 2026 &
Agente autónomo &
Análisis de datos biológicos &
Ejecución autónoma de análisis y hallazgos en notebooks &
La interpretación final sigue requiriendo revisión experta \\
\addlinespace
\textit{Virtual Biotech}, 2026 &
Sistema jerárquico multi-agente &
Descubrimiento terapéutico &
Orquestación científica de punta a punta &
Preprint; evidencia aún preliminar \\
\addlinespace
\textit{Benchmarking LLM-based Agents}, 2026 &
Framework de benchmarking &
Ómicas de célula única &
Evaluación de capacidades agénticas en análisis ómicos &
Benchmarks limitados \\
\addlinespace
\textit{MCPmed}, 2026 &
Arquitectura basada en MCP &
Servicios web bioinformáticos &
Interoperabilidad entre agentes y servicios web bioinformáticos &
Propuesta temprana \\
\addlinespace
\textit{Reliable Data Science Copilots}, 2026 &
Refinamiento iterativo &
Programación científica &
Necesidad de validación iterativa &
Baja precisión inicial sin refinamiento \\
\addlinespace
\textit{Metaviral Workflow Generation}, 2026 &
Generación automática de pipelines &
Virología computacional &
Automatización de workflows multi-step &
Riesgo de alucinaciones \\
\addlinespace
\textit{ggplotAgent}, 2026 &
Agente auto-correctivo &
Visualización científica &
Corrección autónoma de errores &
Alcance limitado \\
\addlinespace
\textit{Functional Genomics Assistant}, 2025 &
Asistente basado en evidencia &
Genómica funcional &
Mejora trazabilidad y procedencia &
Requiere fuentes curadas \\
\addlinespace
\textit{Multi-Agent AI Systems Review}, 2026 &
Survey / taxonomía &
Datos biológicos y clínicos &
Tendencias en arquitecturas multi-agente &
Evidencia mayormente teórica \\
\addlinespace
\textit{AI Agents for Biological Research}, 2026 &
Framework taxonómico 5D &
Automatización científica &
Clasificación estructurada de sistemas agénticos &
Validación experimental limitada \\
\addlinespace
\textit{Automated Protein Protocol Optimization}, 2026 &
Automatización basada en LLMs &
Protocolos de proteínas &
Traducción de literatura a protocolos ejecutables &
Desafíos de integración wet-lab \\
\end{longtable}


% --- SUBSECCIÓN 1 ---
\subsection{Arquitecturas agénticas y patrones de orquestación}\label{subsec:architectures}

Trabajos como \textit{BioMedAgent}~\citep{41912700} y
\textit{PKGPT}~\citep{42076151} mostraron que es posible
integrar herramientas bioinformáticas externas dentro de arquitecturas capaces
de ejecutar workflows complejos y refinar sus resultados de manera iterativa,
aunque todavía bajo supervisión humana en tareas de alta exigencia técnica. De
forma similar, \textit{MCPmed}~\citep{41729821} propuso mecanismos de
interoperabilidad entre agentes y servicios web bioinformáticos mediante
protocolos estructurados, destacando la necesidad de infraestructuras
estandarizadas para sistemas científicos autónomos.

Asimismo, varios estudios identificaron patrones arquitectónicos comunes,
incluyendo sistemas jerárquicos multi-agente, agentes recursivos de refinamiento
y modelos orientados a la orquestación de tareas científicas complejas.

Es importante precisar que esta caracterización arquitectónica no se basa
únicamente en los artículos empíricos incluidos, sino también en una capa
teórica de apoyo sobre diseño de agentes, memoria y uso de herramientas. Por
ello, aunque los trabajos centrados exclusivamente en \textit{prompting} simple
fueron excluidos del corpus principal, el nivel de LLM base se conserva aquí
como referencia conceptual para mostrar el salto estructural hacia sistemas con
planificación, estado y ejecución externa. En ese puente conceptual, la
documentación oficial de LangChain resulta útil para describir de manera
estandarizada cómo se articulan los componentes mínimos de un agente moderno,
incluyendo el ciclo modelo-herramientas, la memoria operativa y la interacción
con funciones externas~\citep{LangChainAgents2026,LangChainMemory2026,LangChainTools2026}.

Esta evolución arquitectónica se puede categorizar en tres estadios
(\cref{fig:agentic-levels}). Inicialmente, los sistemas operaban mediante
\textit{prompting} directo (Nivel 1). La introducción de frameworks como
LangChain permitió el desarrollo de agentes con capacidad de razonamiento y uso
de herramientas (Nivel 2), donde el LLM actúa como un motor de decisiones para
ejecutar software bioinformático. Finalmente, la frontera actual se desplaza
hacia sistemas multi-agente jerárquicos (Nivel 3), donde un supervisor
coordina agentes especializados y distribuye subtareas científicas complejas,
como se observa en propuestas recientes como \textit{Virtual
Biotech}~\citep{41808990} \textit{(preprint)} y en marcos de interoperabilidad
como \textit{MCPmed}~\citep{41729821}.

\begin{figure}[htbp]
  \centering
  \begin{subfigure}[t]{0.47\textwidth}
    \centering
    \includegraphics[width=\linewidth]{llm_base.png}
    \caption{LLM base}
    \label{fig:llm-base}
  \end{subfigure}\hfill
  \begin{subfigure}[t]{0.47\textwidth}
    \centering
    \includegraphics[width=\linewidth]{agente.png}
    \caption{Agente con herramientas}
    \label{fig:agent-tool-use}
  \end{subfigure}

  \vspace{0.8em}

  \begin{subfigure}[t]{0.72\textwidth}
    \centering
    \includegraphics[width=\linewidth]{arquitectura_agentica.png}
    \caption{Arquitectura agéntica}
    \label{fig:agentic-architecture}
  \end{subfigure}
  \caption{Evolución conceptual desde modelos LLM usados mediante
  \textit{prompting} directo hacia agentes con herramientas y arquitecturas
  agénticas capaces de planificar, ejecutar, validar y refinar workflows
  científicos.}
  \label{fig:agentic-levels}
\end{figure}

% --- CAJA DE INFORMACIÓN (INFOBOX) ---
\begin{infobox}[Recuadro 1 \textbar{} Concepto clave]
  Un sistema agéntico combina razonamiento, planificación, memoria operativa,
  uso de herramientas externas y ciclos de retroalimentación. Esta arquitectura
  permite que el LLM deje de funcionar como un generador aislado de texto y pase
  a actuar como coordinador de un workflow computacional verificable.
  
  En bioinformática, esta separación entre razonamiento y ejecución es crítica
  porque permite delegar tareas técnicas a software especializado mientras se
  mantiene una capa de supervisión y validación.
\end{infobox}


Ejemplo de workflow en Python:

\begin{lstlisting}[style=python, caption={Workflow en Python con una libreria
  de LLMs, como LangChain u OpenAI API.}, label=lst:python_example]
from langchain.llms import OpenAI
from langchain.agents import initialize_agent, Tool

# Inicializar LLM
llm = OpenAI(temperature=0.0, model="gpt-4")

# Definir herramientas bioinformaticas
tools = [
    Tool(
        name="PubMedSearch",
        func=search_pubmed,
        description="Buscar articulos relevantes en PubMed"
    ),
    Tool(
        name="ParseSequence",
        func=parse_fasta,
        description="Procesar archivos de secuencias FASTA"
    ),
]

# Crear agente
agent = initialize_agent(tools, llm, agent="zero-shot-react-description")
result = agent.run("Buscar articulos sobre aplicaciones de CRISPR")
\end{lstlisting}

% --- CAJA DE PUNTO CLAVE (KEYPOINT) ---
\begin{keypoint}[Hallazgo clave del eje 1]
  La transición desde prompting directo hacia arquitecturas agénticas permite
  que los LLMs coordinen herramientas externas, ejecuten procesos iterativos y
  estructuren mejor el workflow científico, aunque la validación humana sigue
  siendo necesaria en tareas de alta sensibilidad biológica o clínica.
\end{keypoint}


% --- SUBSECCIÓN 2 ---
\subsection{Automatización de pipelines bioinformáticos específicos}\label{subsec:pipeline-automation}

Una parte significativa de los estudios seleccionados se enfocó en la
automatización de workflows computacionales en áreas como transcriptómica,
\textit{single-cell omics}, modelado farmacocinético y generación automatizada
de protocolos experimentales.

Sistemas como \textit{CellVoyager}~\citep{41845065} demostraron la capacidad de
los agentes para proponer análisis, ejecutarlos y producir resultados de forma
autónoma dentro de notebooks computacionales, especialmente en escenarios de
\textit{single-cell omics}. Otros trabajos mostraron la generación automática
de pipelines multi-step para virología y workflows de análisis ómicos,
reduciendo la intervención manual en tareas repetitivas y de ensamblaje
técnico~\citep{41523654}.

Además, algunos estudios extendieron estas capacidades hacia entornos wet-lab,
traduciendo información contenida en artículos científicos en protocolos
experimentales ejecutables.

\begin{figure}[htbp]
  \centering
  
\includegraphics[width=0.88\textwidth]{workflow.png}
  \caption{Representación conceptual de un workflow bioinformático automatizado
  mediante sistemas agénticos.}
  \label{fig:automated-workflow}
\end{figure}

Un aspecto crítico en esta automatización es la transición hacia workflows de
ciclo cerrado (\textit{closed-loop}). A diferencia de los scripts estáticos,
agentes como \textit{CellVoyager}~\citep{41845065} operan en entornos de ejecución dinámica,
como Jupyter Notebooks, donde el éxito de una tarea depende de la
retroalimentación inmediata del intérprete de código. Esta capacidad de
autocorrección amplía el alcance de procesos complejos, como el análisis de
\textit{single-cell omics}, y desplaza parte de la carga técnica desde el
investigador hacia el sistema, aunque no elimina por completo la necesidad de
revisión experta sobre la interpretación final.

% --- TABLA COMPARATIVA (TABULARX) ---
\begin{table}[H]
  \centering
  \small
  \caption{Comparación de sistemas, herramientas o enfoques evaluados en los estudios.}
  \label{tab:comparison}
  \begin{tabularx}{\textwidth}{@{} l l Y Y @{}}
    \toprule
    \textbf{Sistema} & \textbf{Enfoque} & \textbf{Dominio} &
    \textbf{Contribución principal} \\
    \midrule
    CellVoyager & Agente autónomo & \textit{Single-cell omics} &
    Ejecución autónoma de análisis y generación de hallazgos en notebooks. \\
    PKGPT & Agente recursivo & Farmacometría &
    Refinamiento iterativo de modelos NONMEM mediante ejecución externa y evaluación repetida. \\
    ChatSpatial & Orquestación por esquemas & Transcriptómica espacial &
    Reproducibilidad de workflows entre entornos R y Python mediante esquemas validados. \\
    Alvessa & Asistente con \textit{grounding} & Genómica funcional &
    Verificación a nivel de sentencia y trazabilidad de evidencia. \\
    \bottomrule
  \end{tabularx}
\end{table}

% --- SUBSECCIÓN 3 ---
\subsection{Reproducibilidad, benchmarking y evaluación comparativa}\label{subsec:reproducibility}

La reproducibilidad fue identificada como uno de los principales desafíos
asociados al uso de sistemas basados en agentes dentro de bioinformática.
Diversos estudios propusieron mecanismos para mejorar la trazabilidad,
consistencia y verificabilidad de los workflows automatizados.

\textit{ChatSpatial}~\citep{41890081} introdujo estrategias de orquestación basadas en esquemas
estructurados para mejorar la consistencia entre entornos R y Python, mientras
que otros trabajos propusieron asistentes basados en evidencia y mecanismos de
\textit{grounding} para reducir alucinaciones y mejorar la procedencia de la
información~\citep{41502944}.

De igual manera, algunos estudios desarrollaron frameworks de benchmarking
orientados a evaluar la calidad de generación de código, la ejecución de tareas
analíticas y la robustez de distintas capacidades agénticas en workflows
bioinformáticos complejos.

\begin{figure}[!t]
  \centering
  \includegraphics[width=0.80\textwidth]{observabilidad.png}
  \caption{Observabilidad de sistemas agénticos para auditoría de trazas,
  grafos de ejecución y consumo de recursos en workflows automatizados.}
  \label{fig:observability}
\end{figure}

Finalmente, la adopción de herramientas de observabilidad como
LangSmith~\citep{LangSmithObservability2026,LangSmithConcepts2026} puede ser
especialmente útil para fortalecer la robustez de estos sistemas. Estas
plataformas permiten una auditoría granular de los grafos de ejecución,
facilitando la identificación de cuellos de botella y errores de razonamiento
en flujos multi-agente, lo que eleva el estándar de reproducibilidad al permitir
que investigadores externos inspeccionen cada paso lógico y el consumo de
recursos de los experimentos automatizados.


% --- SUBSECCIÓN 4 ---
\subsection{Limitaciones técnicas y consideraciones de gobernanza}\label{subsec:governance}

A pesar del potencial observado en los sistemas basados en LLMs y agentes, la
literatura revisada reportó múltiples limitaciones relacionadas con
alucinaciones, coste computacional, dependencia de herramientas externas y
problemas de reproducibilidad.

Varios estudios señalaron que los LLMs presentan tasas de error significativas
durante ejecuciones iniciales, haciendo necesaria la incorporación de mecanismos
de validación iterativa, autocorrección y supervisión humana. Asimismo, el
elevado consumo de recursos computacionales y coste asociado al uso continuo de
modelos avanzados representa una barrera importante para su adopción
generalizada.

Finalmente, diferentes autores enfatizaron la ausencia de estándares claros para
la gobernanza y evaluación de sistemas agénticos en entornos científicos,
destacando la necesidad de desarrollar marcos normativos, estrategias de
supervisión humana y protocolos de validación reproducibles para futuras
implementaciones en bioinformática.

\begin{figure}[htbp]
  \centering
  \includegraphics[width=0.88\textwidth]{riesgos.png}
  \caption{La supervisión humana y la observabilidad son componentes críticos
  del workflow agéntico: permiten intervenir en puntos de validación, auditar
  decisiones intermedias y mantener el control experimental durante la
  ejecución de pipelines bioinformáticos automatizados.}
  \label{fig:risks-governance}
\end{figure}

\begin{table}[htbp]
  \centering
  \caption{Limitaciones técnicas y estrategias de mitigación para sistemas
  agénticos en bioinformática.}
  \label{tab:limitations-governance}
  \begin{tabularx}{\textwidth}{@{} l X X @{}}
    \toprule
    \textbf{Categoría} & \textbf{Limitación técnica} &
    \textbf{Estrategia de mitigación} \\
    \midrule
    Fiabilidad &
    Alucinaciones en rutas metabólicas o secuencias. &
    \textit{Grounding} y verificación a nivel de sentencia
    (\textit{Alvessa}~\citealp{41502944}). \\
    Arquitectura &
    Errores en cascada en flujos multi-agente. &
    Implementación de \textit{unit testing} y observabilidad
    (LangSmith~\citealp{LangSmithObservability2026}). \\
    Coste &
    Elevado consumo de tokens en ciclos recursivos. &
    Uso de modelos \textit{small-language} (SLMs) para tareas de soporte. \\
    Gobernanza &
    Falta de trazabilidad en decisiones autónomas. &
    Protocolos de \textit{human-in-the-loop} en puntos críticos de validación. \\
    Interoperabilidad &
    Incompatibilidad entre herramientas biológicas. &
    Estandarización mediante el \textit{Model Context Protocol}
    (MCP~\citealp{41729821}). \\
    \bottomrule
  \end{tabularx}
\end{table}

La literatura sugiere que el despliegue de estos sistemas debe estar regido por
un marco de gobernanza clínica y técnica. Un desafío crítico es el ``coste de la
autonomía''; por ejemplo, sistemas que utilizan razonamiento recursivo pueden
consumir entre 15 y 50 veces más tokens que un enfoque \textit{zero-shot}. Para
mitigar esto, se propone la integración de capas de observabilidad que permitan
monitorizar el presupuesto de tokens en tiempo real. Además, la gobernanza debe
incluir protocolos de validación humana obligatoria antes de la transición de un
flujo computacional (\textit{in silico}) a uno experimental (\textit{in vitro}),
asegurando que la IA actúe como un copiloto de alta velocidad y no como un
tomador de decisiones final sin supervisión.
